\startquestion{1}{一}{填写本问研究对象与应直接回答的问题。}
\subsubsection*{1）模型的准备}
先界定研究对象、输入数据及单位，说明与前问或共用资料的接口。介绍关键变量和参数如何取得，再解释需要解决的困难。

\noindent\hspace*{2em}\topic{资料与变量}
说明实际可用信息、处理方式和时间空间支持。无法识别的量保留未知，不用文字补出尚未取得的数据。

\subsubsection*{2）模型的建立}
从本问目标引出数学关系。式~\eqref{eq:format-demo}仅演示公式排版、编号与引用，正式稿应替换成实际模型或删除。
\begin{equation}
\min_{\mathbf{x}\in\mathcal{X}} f(\mathbf{x};\theta).
\label{eq:format-demo}
\end{equation}
式中，$\mathbf{x}$表示变量，$\theta$表示参数，$\mathcal{X}$表示可行域，$f$表示目标函数。此处只说明数学记号，不意味着本题应使用优化模型。

\subsubsection*{3）模型的求解与分析}
说明实际采用的求解步骤、停止条件和必要核查；结果先回答本问，再解释关键关系。表~\ref{tab:q1-demo}演示三线表，不含研究数值。
\begin{table}[htbp]
\centering
\caption{本问结果展示位置（格式示例）}\label{tab:q1-demo}
\fontsize{9.5}{12}\selectfont
\setlength{\tabcolsep}{4pt}
\begin{tabularx}{\textwidth}{p{25mm}X X}
\toprule
项目 & 正式结果 & 对应解释 \\
\midrule
主要答案 & 填写已确认结果 & 说明它如何回答本问 \\
必要核查 & 填写实际核查结果 & 说明支持范围与未解项 \\
\bottomrule
\end{tabularx}
\end{table}
\FloatBarrier
